\begin{enumerate}
\item \textbf{Random Variable} For each random variable, list all possible values of $X$.
\begin{tasks}[after-skip = 40pt, after-item-skip=40pt](2)
\task$X = $ the number of heads when you toss 10 fair coins.
\task $X = $ the difference of two six-sided dice.
\task $X = $ the number of pairs of cards (cards with the same value) when you draw 5 cards from a single standard deck.
\task $X = $ the temperature of water in Celcius.
\end{tasks}
 
\item \textbf{Probability of Random Variable} Let $Y$ be the sum of two dice. Compute the following
\begin{tasks}[after-skip = 50pt, after-item-skip=50pt](2)
\task $P(Y = 7)$
\task $P(Y \geq 8)$
\task $P(Y \mbox{ is even})$
\task $P(Y \mbox{ is divisible by 3})$
\end{tasks}
 
\item \textbf{Expected Value} Find the expected number of the following random variables.
\begin{tasks}[after-skip = 80pt, after-item-skip=80pt](1)
\task $X=$ the higher die when you roll two six-sided dice. (For example, if you roll a 5 and a 2, the higher value is 5.)
\task $X=$ the number of red marbles drawn from a bag with 8 white marbles and 2 red marbles, when you draw 2 marbles without replacement.
\task $X=$ the sum of a d6, a d8, and a d10. (dN is an N-sided die.)
\end{tasks}

\item You are playing a game by drawing a card from a standard deck and replacing it. If the card is a face card (J,Q,K), you win $\$15$. If it is not a face card, you pay that many \$ (A is considered 1). What is the expected value of playing the game?

\vspace{60pt}

\item \textbf{Probability Distribution}
Let $Y$ be the sum of two dice rolls. Draw the bar chart for the probability distribution function $P(Y)$

\vspace{100pt}

\item \textbf{Binomial Distribution} Suppose that about 85\% of graduating students attend their graduation. A group of 22 graduating students is randomly chosen.
\begin{tasks}(1)
\task In words, define the random variable $X$.
\task List the values that $X$ may take on.
\task How many are expected to attend their graduation?
\task Find the probability that 17 or 18 attend. Express it as a formula involving combination.
\task Based on numerical values, would you be surprised if all 22 attended graduation? Justify your answer numerically.
\end{tasks}

\vspace{100pt}

\item \textbf{Geometric Dsitribution}

According to a recent Pew Research poll, 75\% of millenials (people born between 1981 and 1995) have a profile on a social networking site. Let $X =$ the number of millenials you ask until you find a person without a profile on a social networking site.
\begin{tasks}(1)
\task Find the the expected number $E(X)$.
\task Wkhat is the probability that you must ask ten people to find one person without a social networking site?
\task What is the probability that you must ask 20 people to find one person without a social networking site? Express it as a 
\task What is the probability that you must ask at most five people?
\end{tasks}
\end{enumerate}